\section{Cartan-Eilenberg resolutions}
\label{section-cartan-eilenberg}

\noindent
This section can be expanded. The material can be generalized and applied in
more cases. Resolutions need not use injectives and the method also
works in the unbounded case in some situations.

\begin{definition}
\label{definition-cartan-eilenberg}
Let $\mathcal{A}$ be an abelian category.
Let $K^\bullet$ be a bounded below complex.
A {\it Cartan-Eilenberg resolution} of $K^\bullet$
is given by a double complex $I^{\bullet, \bullet}$
and a morphism of complexes $\epsilon : K^\bullet \to I^{\bullet, 0}$
with the following properties:
\begin{enumerate}
\item There exists a $i \ll 0$ such that $I^{p, q} = 0$ for all $p < i$
and all $q$.
\item We have $I^{p, q} = 0$ if $q < 0$.
\item The complex $I^{p, \bullet}$ is an injective resolution of $K^p$.
\item The complex $\Ker(d_1^{p, \bullet})$ is an injective resolution
of $\Ker(d_K^p)$.
\item The complex $\Im(d_1^{p, \bullet})$ is an injective resolution
of $\Im(d_K^p)$.
\item The complex $H^p_I(I^{\bullet, \bullet})$ is an injective resolution
of $H^p(K^\bullet)$.
\end{enumerate}
\end{definition}

\begin{lemma}
\label{lemma-cartan-eilenberg}
Let $\mathcal{A}$ be an abelian category with enough injectives.
Let $K^\bullet$ be a bounded below complex.
There exists a Cartan-Eilenberg resolution of $K^\bullet$.
\end{lemma}

\begin{proof}
Suppose that $K^p = 0$ for $p < n$. Decompose $K^\bullet$ into
short exact sequences as follows: Set $Z^p = \Ker(d^p)$,
$B^p = \Im(d^{p - 1})$, $H^p = Z^p/B^p$, and consider
$$
\begin{matrix}
0 \to Z^n \to K^n \to B^{n + 1} \to 0 \\
0 \to B^{n + 1} \to Z^{n + 1} \to H^{n + 1} \to 0 \\
0 \to Z^{n + 1} \to K^{n + 1} \to B^{n + 2} \to 0 \\
0 \to B^{n + 2} \to Z^{n + 2} \to H^{n + 2} \to 0 \\
\ldots
\end{matrix}
$$
Set $I^{p, q} = 0$ for $p < n$. Inductively we choose
injective resolutions as follows:
\begin{enumerate}
\item Choose an injective resolution $Z^n \to J_Z^{n, \bullet}$
with $J_Z^{n, m} = 0$ for $m < 0$, see
Lemma \ref{lemma-injective-resolutions-exist}.
\item Using Lemma \ref{lemma-injective-resolution-ses} choose injective
resolutions $K^n \to I^{n, \bullet}$, $B^{n + 1} \to J_B^{n + 1, \bullet}$
with $I^{n, m} = 0$ for $m < 0$ and $J_B^{n + 1, m} = 0$ for $m < 0$,
and an exact sequence of complexes
$0 \to J_Z^{n, \bullet} \to I^{n, \bullet} \to J_B^{n + 1, \bullet} \to 0$
compatible with the short exact sequence
$0 \to Z^n \to K^n \to B^{n + 1} \to 0$.
\item Using Lemma \ref{lemma-injective-resolution-ses} choose injective
resolutions $Z^{n + 1} \to J_Z^{n + 1, \bullet}$,
$H^{n + 1} \to J_H^{n + 1, \bullet}$
with $J_Z^{n + 1, m} = 0$ for $m < 0$ and $J_H^{n + 1, m} = 0$ for $m < 0$
and an exact sequence of complexes
$0 \to J_B^{n + 1, \bullet} \to J_Z^{n + 1, \bullet}
\to J_H^{n + 1, \bullet} \to 0$
compatible with the short exact sequence
$0 \to B^{n + 1} \to Z^{n + 1} \to H^{n + 1} \to 0$.
\item Etc.
\end{enumerate}
Taking as maps $d_1^\bullet : I^{p, \bullet} \to I^{p + 1, \bullet}$
the compositions
$I^{p, \bullet} \to J_B^{p + 1, \bullet} \to
J_Z^{p + 1, \bullet} \to I^{p + 1, \bullet}$ everything is clear.
\end{proof}

\begin{lemma}
\label{lemma-two-ss-complex-functor}
Let $F : \mathcal{A} \to \mathcal{B}$ be a left exact functor of
abelian categories.
Let $K^\bullet$ be a bounded below complex of $\mathcal{A}$.
Let $I^{\bullet, \bullet}$ be a Cartan-Eilenberg resolution
for $K^\bullet$. The spectral sequences
$({}'E_r, {}'d_r)_{r \geq 0}$ and $({}''E_r, {}''d_r)_{r \geq 0}$
associated to the double complex $F(I^{\bullet, \bullet})$
satisfy the relations
$$
{}'E_1^{p, q} = R^qF(K^p)
\quad
\text{and}
\quad
{}''E_2^{p, q} = R^pF(H^q(K^\bullet))
$$
Moreover, these spectral sequences are bounded, converge to
$H^*(RF(K^\bullet))$, and the associated induced filtrations on
$H^n(RF(K^\bullet))$ are finite.
\end{lemma}

\begin{proof}
We will use the following remarks without further mention:
\begin{enumerate}
\item As $I^{p, \bullet}$ is an injective resolution of
$K^p$ we see that $RF$ is defined at $K^p[0]$
with value $F(I^{p, \bullet})$.
\item As $H^p_I(I^{\bullet, \bullet})$ is an injective resolution of
$H^p(K^\bullet)$ the derived functor $RF$ is defined at $H^p(K^\bullet)[0]$
with value $F(H^p_I(I^{\bullet, \bullet}))$.
\item By
Homology, Lemma \ref{homology-lemma-double-complex-gives-resolution}
the total complex $\text{Tot}(I^{\bullet, \bullet})$
is an injective resolution of
$K^\bullet$. Hence $RF$ is defined at $K^\bullet$ with value
$F(\text{Tot}(I^{\bullet, \bullet}))$.
\end{enumerate}
Consider the two spectral sequences associated to the double complex
$L^{\bullet, \bullet} = F(I^{\bullet, \bullet})$, see
Homology, Lemma \ref{homology-lemma-ss-double-complex}.
These are both bounded, converge to $H^*(\text{Tot}(L^{\bullet, \bullet}))$,
and induce finite filtrations on $H^n(\text{Tot}(L^{\bullet, \bullet}))$, see
Homology, Lemma \ref{homology-lemma-first-quadrant-ss}.
Since
$\text{Tot}(L^{\bullet, \bullet}) =
\text{Tot}(F(I^{\bullet, \bullet})) =
F(\text{Tot}(I^{\bullet, \bullet}))$ computes
$RF(K^\bullet)$ we find the final assertion of the lemma holds true.

\medskip\noindent
Computation of the first spectral sequence. We have
${}'E_1^{p, q} = H^q(L^{p, \bullet})$ in other words
$$
{}'E_1^{p, q} = H^q(F(I^{p, \bullet})) = R^qF(K^p)
$$
as desired. Observe for later use that the maps
${}'d_1^{p, q} : {}'E_1^{p, q} \to {}'E_1^{p + 1, q}$ are the maps
$R^qF(K^p) \to R^qF(K^{p + 1})$ induced by $K^p \to K^{p + 1}$
and the fact that $R^qF$ is a functor.

\medskip\noindent
Computation of the second spectral sequence. We have
${}''E_1^{p, q} = H^q(L^{\bullet, p}) = H^q(F(I^{\bullet, p}))$.
Note that the complex $I^{\bullet, p}$ is bounded below,
consists of injectives, and moreover each kernel, image, and
cohomology group of the differentials is an injective object
of $\mathcal{A}$. Hence we can split the differentials, i.e.,
each differential is a split surjection onto a direct summand.
It follows that the same is true after applying $F$. Hence
${}''E_1^{p, q} = F(H^q(I^{\bullet, p})) = F(H^q_I(I^{\bullet, p}))$.
The differentials on this are $(-1)^q$ times $F$ applied to
the differential of the complex $H^p_I(I^{\bullet, \bullet})$
which is an injective resolution of $H^p(K^\bullet)$. Hence the
description of the $E_2$ terms.
\end{proof}

\begin{remark}
\label{remark-functorial-ss}
The spectral sequences of Lemma \ref{lemma-two-ss-complex-functor}
are functorial in the complex $K^\bullet$. This follows from functoriality
properties of Cartan-Eilenberg resolutions. On the other hand, they are
both examples of a more general spectral sequence which may be associated
to a filtered complex of $\mathcal{A}$. The functoriality will follow from
its construction. We will return to this in the section on the filtered
derived category, see Remark \ref{remark-final-functorial}.
\end{remark}










